%% toc: P2 14 14.5,14.6,14.8
\section{Finite-input decision procedures and the arbitrary-language scope}
\label{sec:p02-w02-p02-017-finite-input-scope}

\emph{Purpose.}  Two different kinds of result about classical
propositional logic sit side by side in this chapter.  The decision
procedures of \cref{thm:p02-decidability} and
\cref{def:p02-sat-taut} are algorithms: they read finite binary strings.
The compactness and completeness theorems for an arbitrary set of atoms
(\cref{thm:p02-arbitrary-prop-compactness-bpi,thm:p02-fbc-completeness,%
thm:p01-bpi-propositional-compactness}) are set-theoretic statements:
they speak of premise sets of any size over a set of atoms that need not
be countable or well-orderable, and over \(\mathsf{ZF}\) they are
equivalent to the Boolean Prime Ideal principle.  This section states
exactly where the first kind of result ends and the second begins.  The
separation of the countable, choice-free compactness theorem
\cref{thm:prop-compactness} from the arbitrary-language principle is
already made in \cref{sec:p01-bpi-compactness-bridge} and
\cref{sec:p02-free-boolean-completeness}; it is used here, not repeated.

\emph{Inputs and conclusions.}  The inputs are a set \(P\) of atoms, the
finitary language \(\operatorname{Form}(P)\) of
\cref{def:p02-arbitrary-prop-language,def:p02-fbc-language}, premise
sets \(\Gamma\subseteq\operatorname{Form}(P)\), and the fixed binary
encoding of \cref{def:p02-sat-taut} for the source language
\(\operatorname{Form}(P_{\N})\), where \(P_{\N}=\{p_n:n\in\N\}\) is the
source alphabet of \cref{sec:prop-syntax-semantics-01}.  The conclusions are these.
\begin{enumerate}
 \item The decision procedures act on finite encoded inputs over coded
       atoms.  Every finite input over an arbitrary set of atoms can be
       coded, and its answer does not depend on the coding
       (\cref{thm:p02-w02-p02-017-finite-reduction}).
 \item A set of atoms that is uncountable or not well-orderable admits
       no coding at all, so its language as a whole is not an input
       domain of any such procedure
       (\cref{thm:p02-w02-p02-017-no-coding}).
 \item Even over the coded countable alphabet, the procedures do not
       decide satisfiability of infinite premise sets: there is a
       uniformly decidable family of premise sets whose satisfiability
       is not decidable
       (\cref{thm:p02-w02-p02-017-infinite-undecidable}).
 \item The procedures do not construct the satisfying valuation that
       compactness asserts.  Over the coded alphabet there is a decidable,
       satisfiable premise set with no computable satisfying valuation
       (\cref{thm:p02-w02-p02-017-no-computable-valuation}).  Over an
       arbitrary set of atoms, the finite checks establish only the
       hypothesis of compactness; the valuation itself, such as the
       ultrafilter valuation \(v_U\) of \cref{thm:p02-fbc-separation},
       comes from \(\mathsf{BPI}\)
       (\cref{thm:p02-w02-p02-017-bpi-boundary}).
\end{enumerate}

\subsection{Coded atoms and finite encoded inputs}
\label{sec:p02-w02-p02-017-coded-inputs}

Work in \(\mathsf{ZF}\) throughout; any further principle is named where
it is used.  Satisfaction, satisfiability and semantic consequence for
\(\operatorname{Form}(P)\) are defined from total valuations
\(v:P\to\{0,1\}\) by the recursive clauses of
\cref{sec:prop-syntax-semantics-02}, reading the source values false and true
as \(0\) and \(1\); write \(\Gamma\models_P\varphi\) as in
\cref{thm:p02-fbc-completeness}.  For a formula or finite set of
formulas \(e\), \(\operatorname{Var}(e)\) is the finite set of atoms
occurring in \(e\).  For \(Q\subseteq P\), \(\operatorname{Form}(Q)\) is
identified with the set of \(\alpha\in\operatorname{Form}(P)\) such that
\(\operatorname{Var}(\alpha)\subseteq Q\).

\begin{hbdefinition}{def:p02-w02-p02-017-coded-input}{Atom codings, encoded inputs and effective premise sets}
Let \(Q\) be a set of atoms.
\begin{enumerate}
 \item An \emph{atom coding} of \(Q\) is an injection \(c:Q\to\N\).
       For \(\alpha\in\operatorname{Form}(Q)\), the formula
       \(\alpha^{c}\in\operatorname{Form}(P_{\N})\) replaces each atom \(q\) by
       \(p_{c(q)}\); formally \(q^{c}=p_{c(q)}\), \(\bottom^{c}=\bottom\),
       \((\neg\alpha)^{c}=\neg\alpha^{c}\) and
       \((\alpha\circ\beta)^{c}=(\alpha^{c}\circ\beta^{c})\) for
       \(\circ\in\{\land,\lor,\to\}\).
 \item For \(\gamma_1,\ldots,\gamma_k,\varphi\in\operatorname{Form}(Q)\),
       a \emph{finite encoded input over \(c\)} is the binary string
       that \cref{def:p02-sat-taut} assigns to the finite list of
       formulas \(\gamma_1^{c},\ldots,\gamma_k^{c},\varphi^{c}\), the
       members of the list being separated by one further fixed
       character.  These strings, and nothing else, are what the
       algorithm of \cref{thm:p02-decidability} reads.
 \item For \(\alpha\in\operatorname{Form}(P_{\N})\), the \emph{code number}
       \(\ulcorner\alpha\urcorner\) is the natural number whose binary
       numeral is the digit \(1\) followed by the string that
       \cref{def:p02-sat-taut} assigns to \(\alpha\).  A set
       \(\Gamma\subseteq\operatorname{Form}(P_{\N})\) is \emph{decidable} when
       \(\{\ulcorner\gamma\urcorner:\gamma\in\Gamma\}\) is a decidable
       subset of \(\N\) (\cref{sec:turing-machines-03}).  A family
       \((\Gamma_x)_{x\in\N}\) is \emph{uniformly decidable} when the
       relation
       \(\{(x,\ulcorner\gamma\urcorner):x\in\N,\ \gamma\in\Gamma_x\}\)
       is decidable (\cref{sec:turing-machines-03}).
 \item A valuation \(v:P_{\N}\to\{0,1\}\) is \emph{computable} when
       the function \(n\mapsto v(p_n)\) is computable
       (\cref{def:partial-computable}).
\end{enumerate}
The encoding of \cref{def:p02-sat-taut} is injective, being fully
parenthesized with fixed-width character blocks, and the leading digit
\(1\) keeps leading zeros of the string; hence
\(\alpha\mapsto\ulcorner\alpha\urcorner\) is injective.  A finite set of atoms always has an atom coding, since a
finite set is in bijection with some \(\{0,\ldots,n-1\}\); an
arbitrary set of atoms need not have one.
\index{atom coding}\index{finite encoded input}
\end{hbdefinition}

\begin{hblemma}{thm:p02-w02-p02-017-finite-reduction}{proved=annex:p02-w02-p02-017-finite-reduction}{Finite inputs over an arbitrary set of atoms}
Let \(P\) be any set, \(\Gamma_0\subseteq\operatorname{Form}(P)\)
finite, \(\varphi\in\operatorname{Form}(P)\), and
\(A=\operatorname{Var}(\Gamma_0\cup\{\varphi\})\).
\begin{enumerate}
 \item If valuations \(v,w\) on \(P\) agree on
       \(\operatorname{Var}(\alpha)\), then
       \(\widehat v(\alpha)=\widehat w(\alpha)\).
 \item For \(s:A\to\{0,1\}\), let \(s^{0}\) be the valuation on \(P\)
       equal to \(s\) on \(A\) and to \(0\) off \(A\).  Then
       \(\Gamma_0\) is satisfiable iff \(s^{0}\) satisfies \(\Gamma_0\)
       for one of the \(2^{\abs A}\) functions \(s:A\to\{0,1\}\), and
       \(\Gamma_0\models_P\varphi\) iff every such \(s^{0}\) satisfying
       \(\Gamma_0\) satisfies \(\varphi\).
 \item Let \(Q\subseteq P\) and let \(c\) be an atom coding of \(Q\).
       For every \(\alpha\in\operatorname{Form}(Q)\), every valuation
       \(v\) on \(P\) and every valuation \(u\) on \(P_{\N}\): if
       \(u(p_{c(q)})=v(q)\) for all \(q\in Q\), then
       \(\widehat u(\alpha^{c})=\widehat v(\alpha)\).  Each valuation on
       \(P\) has such a partner on \(P_{\N}\), and each valuation on
       \(P_{\N}\) has such a partner on \(P\), both defined without any
       choice.  Hence, for
       \(\Delta\cup\{\psi\}\subseteq\operatorname{Form}(Q)\),
       \(\Delta\) is satisfiable iff \(\Delta^{c}\) is, and
       \(\Delta\models_P\psi\) iff \(\Delta^{c}\models_{P_{\N}}\psi^{c}\).
 \item Consequently, for every atom coding \(c\) of \(A\), the answers
       of \cref{thm:p02-decidability} on the finite encoded input over
       \(c\) are the satisfiability of \(\Gamma_0\) and the truth of
       \(\Gamma_0\models_P\varphi\); in particular they are the same for
       all atom codings of \(A\).
\end{enumerate}
No choice principle and no form of compactness is used.
\end{hblemma}

The lemma is the exact sense in which the decision procedures reach an
arbitrary set of atoms.  The atoms occurring in one finite input are
finitely many and can be numbered; the question being decided is
defined without reference to the numbering, so no numbering has to be
selected coherently across different inputs.  The procedures are never
applied to the language \(\operatorname{Form}(P)\) as a whole, and the
next lemma shows that for many sets \(P\) they could not be.

\begin{hblemma}{thm:p02-w02-p02-017-no-coding}{proved=annex:p02-w02-p02-017-no-coding}{Only countable well-orderable atom sets can be coded}
If a set \(P\) has an atom coding, then \(P\) is well-orderable and
injects into \(\N\).  The same holds if some injection
\(\operatorname{Form}(P)\to\{0,1\}^{*}\) exists.  Hence if \(P\) is not
well-orderable, or admits no injection into \(\N\) (for example
\(P=\mathcal P(\N)\), by \cref{thm:p01-cantor}), then no injective
encoding of \(\operatorname{Form}(P)\) by finite binary strings exists,
and no algorithm has \(\operatorname{Form}(P)\) as its input domain.
\end{hblemma}

\subsection{Infinite premise sets are not decided}
\label{sec:p02-w02-p02-017-infinite-premises}

An infinite premise set is not a finite string.  The most an algorithm
can receive is a finite description of it, for instance a decision
procedure for membership.  The next theorem shows that such access does
not suffice even when the atoms are the coded alphabet \(P_{\N}\), where
compactness itself holds without any choice principle.  Here \(K\) is the
self-application domain set of \cref{sec:universal-reductions-03}.

\begin{hbtheorem}{thm:p02-w02-p02-017-infinite-undecidable}{proved=annex:p02-w02-p02-017-infinite-undecidable}{Satisfiability of decidable infinite premise sets is undecidable}
There is a uniformly decidable family \((\Gamma_x)_{x\in\N}\) of subsets
of \(\operatorname{Form}(P_{\N})\) such that every \(\Gamma_x\) is decidable and
\[
 \Gamma_x\text{ is satisfiable}\quad\Longleftrightarrow\quad x\notin K .
\]
Consequently the set \(\{x\in\N:\Gamma_x\text{ is satisfiable}\}\) is
neither decidable nor computably enumerable.  For each fixed \(x\), every
finite subset of \(\Gamma_x\) is a finite input whose satisfiability
\cref{thm:p02-decidability} decides.
\end{hbtheorem}

\begin{hbtheorem}{thm:p02-w02-p02-017-no-computable-valuation}{proved=annex:p02-w02-p02-017-no-computable-valuation}{A decidable satisfiable premise set with no computable model}
There is a decidable set \(\Delta\subseteq\operatorname{Form}(P_{\N})\) such that
\(\Delta\) is satisfiable, by a valuation defined in \(\mathsf{ZF}\)
without any choice principle, but no computable valuation satisfies
\(\Delta\).
\end{hbtheorem}

\Cref{thm:p02-w02-p02-017-no-computable-valuation} separates two
questions that the countable compactness theorem joins.  Every finite
subset of \(\Delta\) is satisfiable, and each such fact is settled by
the finite search of \cref{thm:p02-w02-p02-017-finite-reduction}(2),
which also exhibits a satisfying assignment of the finitely many atoms
involved.  Those local assignments are finite objects, yet no computable
procedure assembles a valuation of all of \(\Delta\).  The global valuation is a set-theoretic object whose
existence compactness asserts, not the output of the decision procedure.

\subsection{What \texorpdfstring{\(\mathsf{BPI}\)}{BPI} adds over the finite checks}
\label{sec:p02-w02-p02-017-bpi-adds}

\begin{hbtheorem}{thm:p02-w02-p02-017-bpi-boundary}{proved=annex:p02-w02-p02-017-bpi-boundary}{Finite checks supply the hypothesis, not the valuation}
Work in \(\mathsf{ZF}\).
\begin{enumerate}
 \item For every set \(P\) and every
       \(\Gamma\subseteq\operatorname{Form}(P)\), \(\Gamma\) is finitely
       satisfiable iff, for every finite \(\Gamma_0\subseteq\Gamma\), the
       finite check of \cref{thm:p02-w02-p02-017-finite-reduction}(2)
       succeeds on \(\Gamma_0\).
 \item If \(P\) has an atom coding, every finitely satisfiable
       \(\Gamma\subseteq\operatorname{Form}(P)\) is satisfiable, and this
       uses no choice principle.
 \item The scheme \emph{for every set \(P\), every finitely satisfiable
       \(\Gamma\subseteq\operatorname{Form}(P)\) is satisfiable} is
       equivalent to \(\mathsf{BPI}\).
\end{enumerate}
Consequently, a proof in \(\mathsf{ZF}\) that turns the successful finite
checks into a satisfying valuation, for every set of atoms and every
finitely satisfiable premise set, would be a proof of \(\mathsf{BPI}\) in
\(\mathsf{ZF}\).  In the completeness proof of
\cref{thm:p02-fbc-separation}, the countervaluation \(v_U\) is obtained
from an ultrafilter \(U\) whose existence is the \(\mathsf{BPI}\) step;
no decision procedure computes or selects \(U\).
\end{hbtheorem}

The three clauses locate the boundary.  Clause~1 is all that the finite
procedures contribute: they verify, one finite subset at a time, the
hypothesis of compactness.  Clause~2 is the coded case, where the
choice-free enumeration proof of \cref{thm:prop-compactness} carries the
hypothesis to a valuation; by
\cref{thm:p02-w02-p02-017-no-computable-valuation} that valuation still
need not be computable.  Clause~3 is the arbitrary case, where passing
from the hypothesis to a valuation has exactly the strength of
\(\mathsf{BPI}\).

\begin{hblemma}{thm:p02-w02-p02-017-pairs}{proved=annex:p02-w02-p02-017-pairs}{Valuations of the pair-selection premises are choice functions}
Let \(\mathcal A\) be a set of pairwise disjoint two-element sets, let
\(P=\bigcup\mathcal A\) be the set of atoms, tagged as in
\cref{def:p02-fbc-language}, and put
\[
 \Gamma_{\mathcal A}=
 \{(a\lor b)\land\neg(a\land b):
   a,b\in A\text{ for some }A\in\mathcal A,\ a\ne b\}
 \subseteq\operatorname{Form}(P).
\]
In \(\mathsf{ZF}\): every finite subset of \(\Gamma_{\mathcal A}\) is
satisfiable; a valuation satisfies \(\Gamma_{\mathcal A}\) iff it gives
value \(1\) to exactly one element of each \(A\in\mathcal A\); and
sending a satisfying valuation \(v\) to the function
\(A\mapsto\text{the element of }A\text{ with value }1\) is a bijection
from the satisfying valuations of \(\Gamma_{\mathcal A}\) onto the
choice functions on \(\mathcal A\).  Hence \(\mathsf{ZF}+\mathsf{BPI}\)
proves that every set of pairwise disjoint two-element sets has a choice
function.
\end{hblemma}

\begin{hbexample}{def:p02-w02-p02-017-pairs-example}{Pair selection: finite checks against a global valuation}
\emph{A finite instance.}  Let
\(\mathcal A=\{\{a,b\},\{c,d\}\}\) with \(a,b,c,d\) distinct.  Then
\(\Gamma_{\mathcal A}\) has four members, for instance
\((a\lor b)\land\neg(a\land b)\) and
\((b\lor a)\land\neg(b\land a)\), and
\(\operatorname{Var}(\Gamma_{\mathcal A})=\{a,b,c,d\}\).  The coding
\(a\mapsto0\), \(b\mapsto1\), \(c\mapsto2\), \(d\mapsto3\) turns the first
member into \((p_0\lor p_1)\land\neg(p_0\land p_1)\), and the finite
encoded input is the string of the four coded members.  Of the
\(2^{4}=16\) assignments to \(a,b,c,d\), exactly four satisfy the set:
those making exactly one of \(a,b\) and exactly one of \(c,d\) true.
They correspond to the four choice functions on \(\mathcal A\).  The
coding \(a\mapsto3\), \(b\mapsto2\), \(c\mapsto1\), \(d\mapsto0\) yields
a different input string and the same answer, as
\cref{thm:p02-w02-p02-017-finite-reduction}(4) requires.

\emph{A coded infinite instance.}  Let
\(\mathcal A=\{\{(n,0),(n,1)\}:n\in\N\}\).  The map
\((n,i)\mapsto2n+i\) is an atom coding of \(P=\bigcup\mathcal A\), and
the valuation true exactly on the atoms \((n,0)\) satisfies
\(\Gamma_{\mathcal A}\).  Neither compactness nor \(\mathsf{BPI}\) is
needed, because the family itself names a choice function.

\emph{The uncoded case.}  Suppose \(P=\bigcup\mathcal A\) admits no
injection into \(\N\).  By \cref{thm:p02-w02-p02-017-no-coding} the
language \(\operatorname{Form}(P)\) has no coding.  Each finite subset of
\(\Gamma_{\mathcal A}\) still has a local coding, and every finite check
succeeds by \cref{thm:p02-w02-p02-017-pairs}.  A satisfying valuation of
the whole set, however, is the same thing as a choice function on
\(\mathcal A\).  The finite checks do not produce one; under
\(\mathsf{BPI}\), compactness asserts that one exists.
\end{hbexample}

\begin{hbwarning}{def:p02-w02-p02-017-limits}{What these results do not say}
\begin{enumerate}
 \item The theorems above do not say that \(\mathsf{ZF}\) fails to prove
       \(\mathsf{BPI}\), or fails to prove choice for sets of pairs.  No
       such independence result is proved or cited here;
       \cref{thm:p01-bpi-choice-separation} concerns the gap between
       \(\mathsf{BPI}\) and the Axiom of Choice, not the gap between
       \(\mathsf{ZF}\) and \(\mathsf{BPI}\).  Clause~3 of
       \cref{thm:p02-w02-p02-017-bpi-boundary} is a calibration: whatever
       \(\mathsf{BPI}\) adds to \(\mathsf{ZF}\), the finite checks do not
       supply it.
 \item The answer of a finite check is independent of the coding, but
       the length of the encoded input is not.  It depends on the binary
       lengths of the chosen codes, so complexity bounds for
       \(\mathsf{SAT}\) and \(\mathsf{TAUT}\) in
       \cref{sec:p02-prop-decision-complexity} are statements about
       encoded strings, not about formulas over an uncoded set of atoms.
 \item The failure in \cref{thm:p02-w02-p02-017-infinite-undecidable} is
       a failure for satisfiability.  For the family constructed there,
       unsatisfiability of \(\Gamma_x\) holds exactly when \(x\in K\), a
       computably enumerable condition; the undecidability does not come
       from the absence of compactness, which holds for \(P_{\N}\).
 \item \Cref{thm:p02-w02-p02-017-no-computable-valuation} concerns
       computable valuations of a decidable set over the coded alphabet.
       It does not say that \(\Delta\) lacks a definable valuation; one is
       written down in its proof.
\end{enumerate}
\end{hbwarning}

\emph{Relationships.}  The finite-input procedures are those of
\cref{thm:p02-decidability}, with the encoding and the problems
\(\mathsf{SAT}\) and \(\mathsf{TAUT}\) of \cref{def:p02-sat-taut}.  The
choice-free countable compactness theorem is \cref{thm:prop-compactness};
its relation to the arbitrary-language principle is stated in
\cref{thm:p01-bpi-propositional-compactness} and in the opening of
\cref{sec:p02-free-boolean-completeness}, and the binary-tree route in
\cref{sec:p06-cardinals-choice} is a further countable argument.  The
arbitrary-language statements are
\cref{thm:p02-arbitrary-prop-compactness-bpi,%
thm:p02-arbitrary-prop-completeness,thm:p02-fbc-completeness}, and the
\(\mathsf{BPI}\) step is isolated in \cref{thm:p02-fbc-separation}.  The
computability results used are the numbering and enumeration theorem
\cref{thm:enumeration}, universal computation
\cref{thm:universal-computation}, the normal form
\cref{thm:p04-kleene-normal-form}, the identification of Turing
computability with partial recursiveness
\cref{thm:p04-tm-partial-recursive}, and the undecidability of \(K\) and
non-enumerability of its complement
(\cref{thm:halting-undecidable} and the corollary following it in
\cref{sec:universal-reductions-03}).

\emph{Sources.}  The proofs in the annex are editorial arguments built
only from the handbook results just listed, each with the standing its
own status line records: \cref{thm:p04-kleene-normal-form} and
\cref{thm:p04-tm-partial-recursive} carry cited status, and the form of
the normal form used in the annex, with the minimization defined exactly
when a witness exists, is Kleene's Theorem~XIX(a), equation~(60)
\cite[\S\,63, p.~330]{Kleene1952};
\cref{thm:universal-computation} and the computability of primitive
recursive functions in \cref{sec:turing-machines-02} carry proof
sketches; the remaining results used carry full proofs.  No result in this section is attributed to an
external source beyond those standings.

\index{decision procedure!finite encoded input}
\index{compactness!arbitrary set of atoms}
\index{Boolean Prime Ideal principle!and decision procedures}
